% Options for packages loaded elsewhere
\PassOptionsToPackage{unicode}{hyperref}
\PassOptionsToPackage{hyphens}{url}
\PassOptionsToPackage{dvipsnames,svgnames,x11names}{xcolor}
%
\documentclass[
  letterpaper,
  DIV=11,
  numbers=noendperiod]{scrreprt}

\usepackage{amsmath,amssymb}
\usepackage{iftex}
\ifPDFTeX
  \usepackage[T1]{fontenc}
  \usepackage[utf8]{inputenc}
  \usepackage{textcomp} % provide euro and other symbols
\else % if luatex or xetex
  \usepackage{unicode-math}
  \defaultfontfeatures{Scale=MatchLowercase}
  \defaultfontfeatures[\rmfamily]{Ligatures=TeX,Scale=1}
\fi
\usepackage{lmodern}
\ifPDFTeX\else  
    % xetex/luatex font selection
\fi
% Use upquote if available, for straight quotes in verbatim environments
\IfFileExists{upquote.sty}{\usepackage{upquote}}{}
\IfFileExists{microtype.sty}{% use microtype if available
  \usepackage[]{microtype}
  \UseMicrotypeSet[protrusion]{basicmath} % disable protrusion for tt fonts
}{}
\makeatletter
\@ifundefined{KOMAClassName}{% if non-KOMA class
  \IfFileExists{parskip.sty}{%
    \usepackage{parskip}
  }{% else
    \setlength{\parindent}{0pt}
    \setlength{\parskip}{6pt plus 2pt minus 1pt}}
}{% if KOMA class
  \KOMAoptions{parskip=half}}
\makeatother
\usepackage{xcolor}
\setlength{\emergencystretch}{3em} % prevent overfull lines
\setcounter{secnumdepth}{5}
% Make \paragraph and \subparagraph free-standing
\ifx\paragraph\undefined\else
  \let\oldparagraph\paragraph
  \renewcommand{\paragraph}[1]{\oldparagraph{#1}\mbox{}}
\fi
\ifx\subparagraph\undefined\else
  \let\oldsubparagraph\subparagraph
  \renewcommand{\subparagraph}[1]{\oldsubparagraph{#1}\mbox{}}
\fi


\providecommand{\tightlist}{%
  \setlength{\itemsep}{0pt}\setlength{\parskip}{0pt}}\usepackage{longtable,booktabs,array}
\usepackage{calc} % for calculating minipage widths
% Correct order of tables after \paragraph or \subparagraph
\usepackage{etoolbox}
\makeatletter
\patchcmd\longtable{\par}{\if@noskipsec\mbox{}\fi\par}{}{}
\makeatother
% Allow footnotes in longtable head/foot
\IfFileExists{footnotehyper.sty}{\usepackage{footnotehyper}}{\usepackage{footnote}}
\makesavenoteenv{longtable}
\usepackage{graphicx}
\makeatletter
\def\maxwidth{\ifdim\Gin@nat@width>\linewidth\linewidth\else\Gin@nat@width\fi}
\def\maxheight{\ifdim\Gin@nat@height>\textheight\textheight\else\Gin@nat@height\fi}
\makeatother
% Scale images if necessary, so that they will not overflow the page
% margins by default, and it is still possible to overwrite the defaults
% using explicit options in \includegraphics[width, height, ...]{}
\setkeys{Gin}{width=\maxwidth,height=\maxheight,keepaspectratio}
% Set default figure placement to htbp
\makeatletter
\def\fps@figure{htbp}
\makeatother
% definitions for citeproc citations
\NewDocumentCommand\citeproctext{}{}
\NewDocumentCommand\citeproc{mm}{%
  \begingroup\def\citeproctext{#2}\cite{#1}\endgroup}
\makeatletter
 % allow citations to break across lines
 \let\@cite@ofmt\@firstofone
 % avoid brackets around text for \cite:
 \def\@biblabel#1{}
 \def\@cite#1#2{{#1\if@tempswa , #2\fi}}
\makeatother
\newlength{\cslhangindent}
\setlength{\cslhangindent}{1.5em}
\newlength{\csllabelwidth}
\setlength{\csllabelwidth}{3em}
\newenvironment{CSLReferences}[2] % #1 hanging-indent, #2 entry-spacing
 {\begin{list}{}{%
  \setlength{\itemindent}{0pt}
  \setlength{\leftmargin}{0pt}
  \setlength{\parsep}{0pt}
  % turn on hanging indent if param 1 is 1
  \ifodd #1
   \setlength{\leftmargin}{\cslhangindent}
   \setlength{\itemindent}{-1\cslhangindent}
  \fi
  % set entry spacing
  \setlength{\itemsep}{#2\baselineskip}}}
 {\end{list}}
\usepackage{calc}
\newcommand{\CSLBlock}[1]{\hfill\break\parbox[t]{\linewidth}{\strut\ignorespaces#1\strut}}
\newcommand{\CSLLeftMargin}[1]{\parbox[t]{\csllabelwidth}{\strut#1\strut}}
\newcommand{\CSLRightInline}[1]{\parbox[t]{\linewidth - \csllabelwidth}{\strut#1\strut}}
\newcommand{\CSLIndent}[1]{\hspace{\cslhangindent}#1}

\KOMAoption{captions}{tableheading}
\makeatletter
\@ifpackageloaded{bookmark}{}{\usepackage{bookmark}}
\makeatother
\makeatletter
\@ifpackageloaded{caption}{}{\usepackage{caption}}
\AtBeginDocument{%
\ifdefined\contentsname
  \renewcommand*\contentsname{Table of contents}
\else
  \newcommand\contentsname{Table of contents}
\fi
\ifdefined\listfigurename
  \renewcommand*\listfigurename{List of Figures}
\else
  \newcommand\listfigurename{List of Figures}
\fi
\ifdefined\listtablename
  \renewcommand*\listtablename{List of Tables}
\else
  \newcommand\listtablename{List of Tables}
\fi
\ifdefined\figurename
  \renewcommand*\figurename{Figure}
\else
  \newcommand\figurename{Figure}
\fi
\ifdefined\tablename
  \renewcommand*\tablename{Table}
\else
  \newcommand\tablename{Table}
\fi
}
\@ifpackageloaded{float}{}{\usepackage{float}}
\floatstyle{ruled}
\@ifundefined{c@chapter}{\newfloat{codelisting}{h}{lop}}{\newfloat{codelisting}{h}{lop}[chapter]}
\floatname{codelisting}{Listing}
\newcommand*\listoflistings{\listof{codelisting}{List of Listings}}
\makeatother
\makeatletter
\makeatother
\makeatletter
\@ifpackageloaded{caption}{}{\usepackage{caption}}
\@ifpackageloaded{subcaption}{}{\usepackage{subcaption}}
\makeatother
\ifLuaTeX
  \usepackage{selnolig}  % disable illegal ligatures
\fi
\usepackage{bookmark}

\IfFileExists{xurl.sty}{\usepackage{xurl}}{} % add URL line breaks if available
\urlstyle{same} % disable monospaced font for URLs
\hypersetup{
  pdftitle={tesis\_qmd},
  pdfauthor={Jane Doe},
  colorlinks=true,
  linkcolor={blue},
  filecolor={Maroon},
  citecolor={Blue},
  urlcolor={Blue},
  pdfcreator={LaTeX via pandoc}}

\title{tesis\_qmd}
\author{Jane Doe}
\date{Invalid Date}

\begin{document}
\maketitle

\renewcommand*\contentsname{Table of contents}
{
\hypersetup{linkcolor=}
\setcounter{tocdepth}{2}
\tableofcontents
}
\bookmarksetup{startatroot}

\chapter*{Resumen}\label{resumen}
\addcontentsline{toc}{chapter}{Resumen}

\markboth{Resumen}{Resumen}

Queremos Analizar los datos de la actividad de vigilancia entomológica
de forma oportuna y automática utilizando~herramientas de Ciencia de
Datos utilizando la base de datos de la plataforma de ``Vigilancia
Entomológica y Control Integral del Vector'' del INSTITUTO NACIONAL DE
SALUD PÚBLICA en colaboración con Secretaria de Salud y CENAPRECE. Con
el Fin de calcular y analizar los índices de riesgo entomológico, y con
esto gestionar información para la toma de decisiones en la prevención
de enfermedades trasmitidas por vector.

sistema integral de monitoreo de vectores
http://geosis.mx/aplicaciones/sismv/

ver (Giordano et al. 2020)

Vigilancia Entomológica y Control Integral del Vector
http://geosis.mx/Aplicaciones/EntomologiayControlIntegral/login.aspx

\bookmarksetup{startatroot}

\chapter{Introduccion}\label{introduccion}

\section{Motivacion}\label{motivacion}

Queremos Analizar los datos de la actividad de vigilancia entomológica
de forma oportuna y automática utilizando~herramientas de Ciencia de
Datos utilizando la base de datos de la plataforma de ``Vigilancia
Entomológica y Control Integral del Vector'' del INSTITUTO NACIONAL DE
SALUD PÚBLICA~en colaboración con Secretaria de Salud y el Centro
Nacional de Prevencion y Control de enfermedades (CENAPRECE). Con el Fin
de calcular y analizar los índices de riesgo entomológico, y con esto
gestionar información para la toma de decisiones en la prevención de
enfermedades trasmitidas por vector.

El análisis de riesgo entomológico se hace en hojas de cálculo El manejo
de la información es ineficiente: La estructura de la base de datos es
compleja Adaptar consultas a pequeñas modificaciones es impráctico La
visualización georreferenciada de esta plataforma es pobre e inestable

Usar R para visualizar la información en tiempo y espacio Automatizar el
manejo de la información Generar visualización georreferenciada con los
índices de estegomía

\section{Descripción del problema}\label{descripciuxf3n-del-problema}

El Programa estatal de vigilancia entomología y control integral de
enfermedades trasmitidas por vector, realiza diversas actividades para
la prevención de enfermedades, las cuales se guían con actividades de
vigilancia que proporcionan índices de riesgo entomológico, por ello se
hará una API para calcular los índices de riesgo entomológico de forma
automática y oportuna para la toma de decisione.

Todas las actividades se capturan diariamente en una plataforma, donde
es posible descargar los datos para su reporte y análisis. En el caso de
Vigilancia por Ovitrampa(VO), la plataforma calcula los índices de
riesgo pero para la actividad de Estudio Entomológico (EE) los calcula
por manzana trabajada, estos índices se requiere por localidad de riesgo
de manera semanal.

Se requiere poder calcular índices de riesgo del EE por tipo de estudio,
localidad de riesgo y por semana epidemiológica de manera automatizada.

La actividad de EE evalúa una muestra antes(encuesta) y
después(verificación) de las actividades de control integral para
determinar el riesgo entomológico\\
que hay en un área grande a trabajar (Sectores o Localidad) y si
requiere reforzar acciones de control y prevención.

Tener la información oportuna de los índices de riego entomológicos
permiten tomar decisiones informades de ¿En Dónde? y ¿Cuándo? focalizar
intervenciones.

\section{Gestion de datos}\label{gestion-de-datos}

Para validar y poner a prueba la solución propuesta, se plantea la
utilización de una base de datos ficticia que sea representativa con las
bases de datos reales descargadas de la plataforma. Esto permitirá
evaluar el funcionamiento del paquete de R en un entorno controlado. Una
vez realizadas las pruebas iniciales y asegurada la confiabilidad de los
resultados obtenidos,se establecerá un convenio de confidencialidad con
el fin de acceder a los datos reales y evaluar el desempeño del paquete
de R en un escenario más cercano a la realidad. Esta estrategia
garantizará la confidencialidad de los datos utilizados, a la vez que
brindará la oportunidad de comprobar la efectividad de la solución en
situaciones reales y aplicables al contexto de vigilancia entomológica y
control de enfermedades transmitidas por vectores.

\section{Descripcion de la base de
datos}\label{descripcion-de-la-base-de-datos}

\bookmarksetup{startatroot}

\chapter{Prelimilares}\label{prelimilares}

Guias operaticas de WHO para la vigilancia entomologica

\section{Vigilancia entomologica}\label{vigilancia-entomologica}

Vigilancia de entomologica es importante para determinar la
distribución, la densidad de población, los principales hábitats de las
larvas y los factores de riesgo espaciales y temporales relacionados con
la transmisión de arbovirosis, y los niveles de susceptibilidad o
resistencia a los insecticidas,con el fin de priorizar áreas y
estaciones para el control de vectores. Estos datos permitirán la
selección y el uso de las herramientas de control de vectores más
apropiadas y pueden usarse para monitorear su efectividad. Existen
varios métodos disponibles para la detección y seguimiento de
poblaciones larvarias y adultos. La selección de métodos apropiados
depende de los objetivos de vigilancia, los niveles de infestación y la
disponibilidad de recursos.

\section{Indices estegomia}\label{indices-estegomia}

Las metodologías más comunes emplean procedimientos de muestreo de
larvas en lugar de recolecciones de huevos o adultos. La unidad básica
de muestreo es la casa o localidad, donde se busca sistemáticamente
recipientes para almacenar agua.

\subsection{indice de casa positiva
(HI)}\label{indice-de-casa-positiva-hi}

El índice de casa positiva se ha utilizado más ampliamente para
monitorear los niveles de infestación, pero no toma en cuenta la
cantidad de contenedores positivos ni la productividad de esos
contenedores.

\subsection{indice de contenedor positivo
(CI)}\label{indice-de-contenedor-positivo-ci}

El índice de contenedores sólo proporciona información sobre la
proporción de contenedores con agua que son positivos con larvas o pupas
de mosquito.

Los contenedores se examinan para detectar la presencia de larvas y
pupas de mosquitos. Dependiendo de los objetivos del estudio, la
búsqueda podrá finalizar tan pronto como se encuentren larvas de
mosquito, o podrá continuar hasta que se hayan examinado todos los
contenedores. Es necesaria la recolección de especímenes para examinar
en laboratorio para confirmar las especies presente.

\subsection{indice de bretau (BI)}\label{indice-de-bretau-bi}

El índice Breateau establece una relación entre contenedores positivos y
casas, y se considera el más informativo, pero nuevamente no refleja la
productividad de los contenedores. Sin embargo, en el curso de la
recopilación de información básica para calcular el índice de Breateau,
es posible y deseable obtener un perfil de las características del
hábitat de las larvas registrando simultáneamente la abundancia relativa
de los diversos tipos de contenedores, ya sea como sitios potenciales o
reales de producción de mosquitos. (por ejemplo, número de tambores
positivos por 100 casas, número de neumáticos positivos por 100 casas,
etc.). Estos datos son particularmente relevantes para centrar esfuerzos
para la gestión o eliminación de los hábitats más comunes y para la
orientación de mensajes educativos en apoyo de iniciativas comunitarias.

\section{Definicion y calculo de indices de
estegomia}\label{definicion-y-calculo-de-indices-de-estegomia}

agregar la definision y uso de la vigilancia de las guias de WHO

\bookmarksetup{startatroot}

\chapter{Cálculo de los índices de
estegomía}\label{cuxe1lculo-de-los-uxedndices-de-estegomuxeda}

\section{}\label{section}

Se sigue el modelo de desarrollo basado test driven development Se
desarrollaron basado en tidyr y devtools tidyr Manipulación de los datos
Estilo de codificación devtools Documentación (roxygen2 ) Pruebas
unitarias (testthat) CRAN (Comprehensive R Archive Network)

\bookmarksetup{startatroot}

\chapter{References}\label{references}

\phantomsection\label{refs}
\begin{CSLReferences}{1}{0}
\bibitem[\citeproctext]{ref-GIORDANO2020}
Giordano, Rosanna, Ravi Kiran Donthu, Aleksey V. Zimin, Irene Consuelo
Julca Chavez, Toni Gabaldon, Manuella van Munster, Lawrence Hon, et al.
2020. {``Soybean Aphid Biotype 1 Genome: Insights into the Invasive
Biology and Adaptive Evolution of a Major Agricultural Pest.''}
\emph{Insect Biochemistry and Molecular Biology} 120: 103334.
https://doi.org/\url{https://doi.org/10.1016/j.ibmb.2020.103334}.

\end{CSLReferences}



\end{document}
